\section{第7章 微分方程}
\subsection{微分方程的基本概念}
\begin{exercise}{721}{微分方程辨识与阶数}
设 $y=y(x)$，分别判断 $x^2+x+1=0$、$y^2+2x+1=0$、$x^2y^2+y=3$、$x^2+y'+y=2$ 是否为微分方程；对符合者给出阶数及是否线性的判断，说明未知函数的代数幂次为何不能用来判定阶数。
\end{exercise}
\begin{exercise}{722}{导数阶数与非线性}
判断 $y''-2x(y')^3+y^2=0$ 的阶数及线性性质。给定 $y(0)=a,y'(0)=b$，求该解应满足的 $y''(0)$ 与 $y'''(0)$；解释 $(y')^3$ 的指数为何不能视为导数阶数。
\end{exercise}
\begin{exercise}{723}{微分形式与标准形式}
将 $x\dd y+(x+y)\dd x=0$ 在 $x>0$ 上写为标准形式并判断阶数。求通解及满足 $y(1)=0$ 的解，说明化为标准形式时对 $x=0$ 作了什么限制。
\end{exercise}
\begin{exercise}{724}{特解验证与通解区分}
验证 $y=5x^2$ 是 $xy'-2y=0$ 的解，说明它是通解还是特解。求该方程在 $x>0$ 上的通解及满足 $y(1)=5$ 的解，并解释为何讨论通解时需要指定区间。
\end{exercise}
\begin{exercise}{725}{解族验证与初值确定}
验证 $y=C_1\cos x+C_2\sin x$ 是 $y''+y=0$ 的解，并证明两个常数独立。求满足 $y(0)=2$、$y'(0)=3$ 的特解及其最大值。
\end{exercise}
\begin{exercise}{726}{任意常数个数与通解判定}
验证 $y=x^2/2+x+C$ 满足 $y''-y'=-x$。判断该一参数解族是否为通解；求完整通解，并给出一个不属于上述一参数族的解。说明仅验证代入成立为何还不能判定为通解。
\end{exercise}
\begin{exercise}{727}{隐函数解与正则分支}
在 $x>0$、$y>0$、$y\ne1$ 的正则分支上，验证 $y=\ln(xy)$ 所确定的函数满足
\[
x(y-1)y''+x(y')^2+(y-2)y'=0.
\]
说明该隐式关系给出通解还是特解，并解释为何必须排除 $y=1$ 处直接视作 $y(x)$ 求导。
\end{exercise}
\begin{exercise}{728}{解族、奇解与初值唯一性}
验证 $y=(x-C)^2$ 在 $x\ge C$ 上满足 $y'=2\sqrt y$，并验证 $y\equiv0$ 也是解。对初值 $y(0)=0$，构造至少两个在 $[0,+\infty)$ 上不同的解，说明直接分离变量可能遗漏什么解。
\end{exercise}
\begin{exercise}{729}{消去常数与微分方程阶数}
消去曲线族 $y=C_1e^x+C_2e^{-x}+x$ 中的任意常数，建立二阶微分方程。再求该方程满足 $y(0)=1$、$y'(0)=0$ 的特解，并核验两个初始条件。
\end{exercise}
\subsection{可分离变量的微分方程}
\begin{exercise}{730}{可分离变量方程辨识}
分别判断 $y'+xy=0$、$y^2\dd x+(2x+1)\dd y=0$、$y'=xye^{x+y}$、$(x-y)\dd x+(x+y)\dd y=0$ 是否可直接分离变量。对符合者写出分离形式，并检查分离时被除因子为零可能给出的常值解；说明最后一个方程的右端为何不能表示为单纯的 $x$ 函数与 $y$ 函数之积。
\end{exercise}
\begin{exercise}{731}{变量分离与奇点两侧解}
求 $y'=y(1-x)/x$ 在 $x>0$ 上的通解，包括零解；求满足 $y(1)=e^{-1}$ 的解。讨论该特解能否延拓为在全实轴上满足 $xy'=(1-x)y$ 的可导解。
\end{exercise}
\begin{exercise}{732}{变量分离与解的定义区间}
求 $y'=\sin x\cos x/y^2$ 的隐式通解。求满足 $y(0)=1$ 的解及其在实轴上的有界性；说明为什么通解中必须限制 $y\ne0$，而不能把零函数补作解。
\end{exercise}
\begin{exercise}{733}{分离变量与初值问题}
求 $(1+e^x)\dd y-ye^x\dd x=0$ 满足 $y(0)=2$ 的特解。推广至 $y(0)=a$，判断何种 $a$ 使解恒正、恒负或恒为零，并求 $y(x)/(1+e^x)$。
\end{exercise}
\begin{exercise}{734}{变量分离与平方根分支}
求 $(xy^2+x)\dd x+(y-x^2y)\dd y=0$、$y(0)=1$ 的特解。分离变量时说明需要避开的因子，确定包含 $0$ 的最大可导解区间，并解释不能将正、负平方根在零点任意拼接。
\end{exercise}
\begin{exercise}{735}{隐式解与初值分支}
求 $e^{x-y}\dd x+y\dd y=0$ 满足 $y(0)=1$ 的特解，以隐式形式表示即可。确定包含 $0$ 且使该解保持 $y>0$ 的最大区间，说明区间端点出现的原因。
\end{exercise}
\begin{exercise}{736}{切线斜率与曲线恢复}
曲线通过 $(0,-1)$，其任一点 $(x,y)$ 处切线斜率为 $y^2/(1+x^2)$。求曲线方程及包含 $0$ 的最大存在区间，并求曲线在 $x\to+\infty$ 时的水平渐近线（若该方向属于解的定义域）。
\end{exercise}
\begin{exercise}{737}{积分方程与链式求导}
设连续函数 $f$ 满足 $f(x)=\int_0^{3x}f(t/3)\dd t+2$。证明 $f$ 可导并求 $f$，再直接核验积分等式；若将外部常数 $2$ 改为 $a$，说明解如何改变。
\end{exercise}
\begin{exercise}{738}{含导数积分方程与非线性降阶}
设 $f\in C^1[0,+\infty)$，满足 $f(x)=2\int_0^x(x-t)f(t)f'(t)\dd t+x-1$。先恢复 $f(0)$，再化为一阶初值问题求 $f$；证明所得表达式在整个规定区间上有效。
\end{exercise}
\begin{exercise}{739}{卷积积分方程与初值恢复}
设连续函数 $f$ 满足 $\int_0^xf(t)\dd t=\int_0^xtf(x-t)\dd t+x$。通过换元将等式化为可求导形式，求 $f$，并说明初值从何而来。
\end{exercise}
\begin{exercise}{740}{非局部条件与常数确定}
设非零连续函数 $f$ 满足 $f(x)=\int_0^xf(t)\dd t+2\int_0^1tf^2(t)\dd t$。求 $f$，并证明非零条件下解唯一。若允许零函数，列出全部连续解并代回核验。
\end{exercise}
\begin{exercise}{741}{线性方程与积分条件}
设 $f'=xf/(1+x^2)$、$f(0)=a$，且 $\int_0^{\sqrt3}xf(x)\dd x=7$。求 $a$ 及 $f$，说明积分条件能否唯一确定 $a$；再求 $f$ 在实轴上的最小值。
\end{exercise}
\begin{exercise}{742}{非线性方程的平移代换}
求 $y'=y^2+2(\sin x-1)y+\sin^2x-2\sin x-\cos x+1$ 的通解。采用 $u=y+\sin x-1$ 化简，并补全分离变量可能遗漏的解；求满足 $y(0)=1$ 的解。
\end{exercise}
\begin{exercise}{743}{乘积代换与隐式通解}
求 $y(xy+1)\dd x+x(1+xy+x^2y^2)\dd y=0$ 的隐式通解。采用 $u=xy$，在 $xy\ne0$ 的区间上化为可分离形式并积分；另检查 $y\equiv0$，并通过全微分核验所得隐式关系。
\end{exercise}
\begin{exercise}{744}{平衡解与最大存在区间}
求初值问题 $y'=y^2$、$y(0)=a$ 的解，包括 $a=0$ 的情形；分别确定 $a>0$ 与 $a<0$ 时包含 $0$ 的最大存在区间，并讨论 $x\to+\infty$ 时解是否有意义。
\end{exercise}
\begin{exercise}{745}{受限增长与阈值时间}
种群满足 $N'=rN(1-N/K)$，其中 $r,K>0$、$N(0)=N_0\in(0,K)$。求 $N(t)$，证明其单调增加且趋于 $K$；当 $N_0<K/2$ 时，求首次达到 $K/2$ 的时间。
\end{exercise}
\subsection{齐次方程}
\begin{exercise}{746}{齐次方程辨识与求解}
分别判断 $y'+y=0$、$y^2\dd x+x\dd y=0$、$y'+xy=e^x$、$x^2y'+xy=y^2$ 是否为形如 $y'=\varphi(y/x)$ 的齐次方程。对符合者求 $x>0$ 上的通解，并补全常比例解。
\end{exercise}
\begin{exercise}{747}{根式齐次方程与漏解}
求 $y'=2\sqrt{y/x}+y/x$ 在 $x>0,y>0$ 上的通解，明确表达式成立的区间限制；另验证 $y\equiv0$ 为解。求满足 $y(1)=1$ 的正解及其保持为正的最大区间。
\end{exercise}
\begin{exercise}{748}{三角齐次方程与分支限制}
求 $y'=\sec(y/x)+y/x$ 在 $x>0$ 上的隐式通解。再求满足 $y(1)=0$ 的解在包含 $1$ 的区间上的显式表达式，并确定该区间端点，说明反三角函数分支的选择。
\end{exercise}
\begin{exercise}{749}{齐次初值问题与非零约束}
求 $y'=(x^2+y^2)/(xy)$、$y(1)=2$ 的特解。给出包含 $1$ 的最大存在区间，证明所选分支始终为正，并判断其能否以原方程解的形式跨越左端点。
\end{exercise}
\begin{exercise}{750}{切线斜率模型与有限区间}
曲线通过 $(1,0)$，切线斜率为 $1+y/x+(y/x)^2$。求曲线方程及包含 $x=1$ 的最大存在区间；说明方程的三角函数表达式为何不能跨过相邻极点继续当作同一解。
\end{exercise}
\begin{exercise}{751}{齐次代换与平衡比例解}
求 $y(x-y)\dd x+x^2\dd y=0$ 满足 $y(1)=1$ 的特解及包含 $1$ 的最大存在区间。求通解过程中，列出将 $y/x$ 视为常数时得到的解，并判断初值解是否属于其中之一。
\end{exercise}
\begin{exercise}{752}{对数齐次方程与正值定义域}
求 $xy'+y(\ln x-\ln y)=0$、$y(1)=e^3$ 的特解，明确对数要求的定义域。再将初值改为 $y(1)=a>0$，求解并讨论所得解在 $x>0$ 上的单调性。
\end{exercise}
\begin{exercise}{753}{交换因变量与指数齐次代换}
求 $(1+2e^{x/y})\dd x+2e^{x/y}(1-x/y)\dd y=0$ 的隐式通解，限定 $y\ne0$。利用 $u=x/y$ 或识别全微分完成推导，并通过微分核验结果。
\end{exercise}
\begin{exercise}{754}{齐次代换与三角隐式解}
求 $(2y\sin(x/y)+3x\cos(x/y))\dd y-3y\cos(x/y)\dd x=0$ 在 $y\ne0$ 上的通解。说明分离变量中除以 $\sin(x/y)$ 是否遗漏常比例解，列出所遗漏的情形。
\end{exercise}
\begin{exercise}{755}{齐次代换与不可跨越的奇点}
求微分方程 $xy'=y+x e^{y/x}$ 在 $x>0$ 上满足 $y(1)=0$ 的解，并确定包含 $1$ 的最大存在区间。证明所得解无法在该区间右端连续延拓为有限值。
\end{exercise}
\begin{exercise}{756}{齐次方程与代换遗漏解}
求 $xy'=y+y^2/x$ 在 $x>0$ 上的全部解，包括代换分离时可能遗漏的解；对初值 $y(1)=a$，确定解的表达式及包含 $1$ 的最大存在区间。
\end{exercise}
\subsection{一阶线性微分方程}
\begin{exercise}{757}{齐次线性方程与初值符号}
求 $y'+2xy=0$ 的通解。对初值 $y(0)=a$，求解在全实轴上的有界性、符号及 $x\to\pm\infty$ 时的极限，并讨论 $a=0$。
\end{exercise}
\begin{exercise}{758}{初值问题与系数奇点}
求 $y'=y\tan x$、$y(0)=1$ 的特解。确定包含 $0$ 的最大存在区间，并求该区间内的最小值，说明不能直接把同一表达式跨过 $\pi/2$ 作为初值解。
\end{exercise}
\begin{exercise}{759}{一阶线性方程辨识与积分形式解}
分别判断 $y'+y^2=0$、$(y')^2=\sqrt{xy}$、$y'+xy=e^xy+x^2$、$y'+\sin(xy)=1$ 是否为关于 $y$ 的一阶线性方程。将符合者写为标准形式，给出积分因子及满足 $y(0)=1$ 的积分形式解，不要求化为初等函数。
\end{exercise}
\begin{exercise}{760}{一阶线性方程与平衡解}
求 $y'+2xy=4x$ 的通解。对初值 $y(0)=a$，讨论解与直线 $y=2$ 的位置关系及两端极限；求恰为常数解的 $a$。
\end{exercise}
\begin{exercise}{761}{积分因子与有界性}
求 $y'+y\cos x=e^{-\sin x}$ 的通解。判断是否存在在整个实轴上有界的解，以及是否存在周期解，分别给出依据。
\end{exercise}
\begin{exercise}{762}{奇异系数与原方程延拓}
求 $(x+1)y'-y=e^x(x+1)^2$、$y(0)=2$ 的特解。先在 $x>-1$ 上求解，再检验所得函数能否延拓为全实轴上的原方程解；指出将方程除以 $x+1$ 会失去什么信息。
\end{exercise}
\begin{exercise}{763}{积分因子与对数外力}
求 $y'+y\cos x-e^{-\sin x}\ln x=0$、$y(1)=0$ 的特解。明确最大定义区间，判断 $x=1$ 处是否为极值点，并给出判断依据。
\end{exercise}
\begin{exercise}{764}{交换自因变量的线性化}
求 $(y^2-6x)y'+2y=0$ 的隐式通解。先将 $x$ 视为 $y$ 的函数并在线性化时注明 $y\ne0$，再单独检查 $y\equiv0$ 是否为原方程的解。
\end{exercise}
\begin{exercise}{765}{关于逆函数的一阶线性方程}
求 $y^3\dd x+(2xy^2-1)\dd y=0$ 的隐式通解。说明将 $x$ 看作 $y$ 的函数时的定义域限制，并补全可能遗漏的常值解；给出通过 $(0,1)$ 的局部分支。
\end{exercise}
\begin{exercise}{766}{隐式初值解与垂直切线}
求 $y\dd x+(x-3y^2)\dd y=0$、$y(1)=1$ 的特解，以隐式形式表示。求初始点附近 $y'(x)$，并确定该分支首次出现垂直切线时的点。
\end{exercise}
\begin{exercise}{767}{逆函数线性化与局部可解性}
求 $y'=1/(x+y)$ 的隐式通解。以 $x$ 为 $y$ 的函数求解，说明在什么点可局部恢复为 $y(x)$；再求通过 $(0,1)$ 的隐式特解及该点导数。
\end{exercise}
\begin{exercise}{768}{线性方程与全微分核验}
求 $y\dd x+(x+y)\dd y=0$ 的隐式通解，并检查 $y\equiv0$。对通过 $(1,1)$ 的分支，求 $y'(1)$，说明隐式解在该点附近能表示为函数的依据。
\end{exercise}
\begin{exercise}{769}{切线斜率与积分因子}
曲线 $y=f(x)$ 通过 $(0,1)$，任一点切线斜率为 $\sec x-f(x)\tan x$。求曲线方程及包含 $0$ 的最大存在区间，计算 $f'(0)$ 并判断该点是否为极值点。
\end{exercise}
\begin{exercise}{770}{积分方程与线性初值问题}
设连续函数 $f$ 满足 $f(x)=\int_0^xf(t)\dd t-x^2$。求 $f$，说明连续性为何足以保证所需可导性；再代回核验，并求 $f$ 在原点处最低阶非零泰勒项。
\end{exercise}
\begin{exercise}{771}{卷积换元与二次恢复}
设连续函数 $f$ 满足 $\int_0^xf(x-t)\dd t=\int_0^x(x-t)f(t)\dd t+e^{-x}-1$。通过换元和求导求 $f$，明确求导后获得的初值，并核验原积分等式。
\end{exercise}
\begin{exercise}{772}{积分因子与有界解}
求方程 $y'+2y=1+e^{-x}$ 的通解；判断是否存在在整个实轴上有界的解，证明你的判断。再求所有解在 $x\to+\infty$ 时的极限。
\end{exercise}
\begin{exercise}{773}{分段外力与连续拼接}
某系统满足 $y'+y=u(x)$、$y(0)=0$，其中 $u(x)=1$（$0\le x<1$）、$u(x)=0$（$x>1$）。分别求两个区间上的解，使 $y$ 在 $x=1$ 连续；判断 $y$ 在 $x=1$ 是否可导，并说明理由。
\end{exercise}
\subsection{可降阶的高阶微分方程}
\begin{exercise}{774}{连续积分与初值数目}
求 $y^{(4)}=x$ 的通解，说明四个积分常数独立。再求满足 $y(0)=y'(0)=y''(0)=0$、$y'''(0)=1$ 的特解，并核验四个条件。
\end{exercise}
\begin{exercise}{775}{连续积分与三角指数外力}
求 $y''=e^{3x}+\sin2x$ 的通解，再求满足 $y(0)=y'(0)=0$ 的解。分别对所得表达式求二阶导数，核验两种外力的积分系数与符号。
\end{exercise}
\begin{exercise}{776}{连续积分与指数多项式}
求 $y''=xe^x$ 的通解及满足 $y(0)=y'(0)=0$ 的特解。求该特解与 $xe^x$ 在 $x\to+\infty$ 时的比值极限，并说明积分常数是否影响该极限。
\end{exercise}
\begin{exercise}{777}{反三角原函数与初值性质}
求 $y''=1/(1+x^2)$ 的通解。再求 $y(0)=y'(0)=0$ 的解，证明它是偶函数、恒非负且严格凸，并确定唯一最小值点。
\end{exercise}
\begin{exercise}{778}{三阶积分与三个初值}
求 $y'''=e^{2x}$ 满足 $y(1)=y'(1)=y''(1)=0$ 的特解。将结果写成 $e^{2x}$ 减去其在 $x=1$ 的二次泰勒多项式的形式并给出系数，再证明解在 $x>1$ 时为正。
\end{exercise}
\begin{exercise}{779}{切线条件与曲线恢复}
曲线经过 $(0,1)$，在该点与 $y=x/2+1$ 相切，并满足 $y''=x$。求曲线方程，证明两项几何条件分别对应两个独立初值，并判断 $(0,1)$ 是否为拐点。
\end{exercise}
\begin{exercise}{780}{不显含因变量的降阶}
求 $y''=y'+x$ 的通解。确定初始斜率 $y'(0)$ 应取何值，才能使解不含 $e^x$ 项；在此条件下，再由 $y(0)=1$ 求唯一解并核验。
\end{exercise}
\begin{exercise}{781}{导数代换与三角分支}
求 $y''+(y')^2+1=0$ 的通解。再求满足 $y(0)=y'(0)=0$ 的特解及包含 $0$ 的最大存在区间，说明积分得到的对数绝对值与区间选择之间的关系。
\end{exercise}
\begin{exercise}{782}{奇异系数与端点有界性}
求 $xy''+3y'=0$ 在 $x>0$ 上的通解及满足 $y(1)=0,y'(1)=2$ 的特解。确定通解中在 $x\to0^+$ 时有界的全部解，说明化为标准形式时不能包含 $x=0$。
\end{exercise}
\begin{exercise}{783}{非线性速度方程与有限奇点}
求 $y''-2(y')^2=0$、$y(0)=0,y'(0)=-1$ 的特解。令 $p=y'$ 降阶并求包含 $0$ 的最大存在区间，判断能否在左端点连续延拓为有限值，并验证原初值。
\end{exercise}
\begin{exercise}{784}{速度线性方程与单调性}
求 $(1+x^2)y''=2xy'$、$y(0)=1,y'(0)=3$ 的特解。给出降阶后速度方程的积分因子，判断解在全实轴上的单调性及零点个数，并用方程核验。
\end{exercise}
\begin{exercise}{785}{不显含自变量与常值解}
求 $yy''+2(y')^2=0$ 在 $y\ne0$ 区间上的通解，并检查常值解是否包含其中。采用 $p=y'$ 降阶，或利用 $y^3$ 的二阶导数求解；说明非常值解为何不能以可导方式跨越其零点。
\end{exercise}
\begin{exercise}{786}{乘积求导与正值初始分支}
求 $yy''+(y')^2=0$、$y(0)=1,y'(0)=1/2$ 的特解。由 $(yy')'=0$ 降阶，确定包含 $0$ 的最大存在区间，并说明平方后求解时如何选取符号。
\end{exercise}
\begin{exercise}{787}{不显含自变量的降阶与能量式}
求 $y^3y''=1$ 在 $y\ne0$ 区间上的隐式通解，先令 $p=y'$ 推导一阶积分。再求 $y(0)=1$、$y'(0)=0$ 的特解，确定其在全实轴上的正负与最小值。
\end{exercise}
\begin{exercise}{788}{不显含因变量的降阶与初值}
求 $y''+(y')^2=1$、$y(0)=y'(0)=0$ 的特解。令 $p=y'$ 降阶后，分别验证求得的速度和原函数；讨论解在全实轴上的奇偶性与凹凸性。
\end{exercise}
\begin{exercise}{789}{不显含因变量的降阶}
求初值问题 $xy''+y'=x$、$y(1)=0$、$y'(1)=1$ 在 $x>0$ 上的解，判断该解能否以二阶可导方式延拓到 $x=0$，并给出依据。
\end{exercise}
\begin{exercise}{790}{不显含自变量与能量积分}
设 $y''=-y^3$、$y(0)=a>0$、$y'(0)=0$。通过降阶证明 $\frac12(y')^2+\frac14y^4=\frac14a^4$；在解首次降到 $0$ 之前确定 $y'$ 的符号，并用定积分表示首次到达 $0$ 的时间，证明该时间有限。
\end{exercise}
\subsection{高阶线性微分方程}
\begin{exercise}{791}{函数组线性相关性}
在实轴上分别判断 $\{x^2,2x^2\}$、$\{\sin x,\cos x\}$、$\{\sin2x,\sin x\cos x\}$、$\{e^{2x},3e^{2x}\}$ 的线性相关性。对无关者给出证明；说明将定义区间缩小为任意非空开区间是否改变结论。
\end{exercise}
\begin{exercise}{792}{已知解验证与基本解组}
对方程 $(x-1)y''-xy'+y=0$，分别判断 $\{x,2x\}$、$\{e^x,8e^x\}$、$\{x,x^2\}$、$\{x,e^x\}$ 能否构成 $x>1$ 上的基本解组。求通解，并说明为何需要把“两个解”与“两个线性无关的解”分开检验。
\end{exercise}
\begin{exercise}{793}{一阶非齐次方程的解结构}
设 $p,f$ 在区间 $I$ 上连续，$y_1$ 是 $y'+py=f$ 的一个解，$y_2\ne0$ 是对应齐次方程的解。证明全部解恰为 $y=y_1+Cy_2$，并用初值 $y(x_0)=a$ 表示 $C$。说明若误把 $C$ 乘在 $y_1$ 前会产生什么条件限制。
\end{exercise}
\begin{exercise}{794}{一阶齐次解与冗余常数}
设连续系数方程 $y'+p(x)y=0$ 在同一区间上有两个不同解 $y_1,y_2$。证明 $y=C(y_1-y_2)$ 为通解。判断 $C_1y_1+C_2y_2$ 中两常数是否独立，并讨论其中一个已知解为零函数的情形。
\end{exercise}
\begin{exercise}{795}{朗斯基行列式与通解判据}
设 $f_1,f_2$ 是 $y''+py'+qy=0$ 的两个非零解，$p,q$ 为常数。记 $W=f_1f_2'-f_2f_1'$。证明若某点 $W\ne0$，则 $C_1f_1+C_2f_2$ 是通解；推导 $W'=-pW$，并说明为什么“两个不同的非零解”本身还不足够。
\end{exercise}
\begin{exercise}{796}{三角基本解组与退化参数}
验证 $\cos ax,\sin ax$ 是 $y''+a^2y=0$（$a\ne0$）的线性无关解并写出通解。若 $a=0$，说明该函数组为何失效并另写通解；求两种情况下满足 $y(0)=0,y'(0)=1$ 的解。
\end{exercise}
\begin{exercise}{797}{非齐次通解的完整性验证}
验证 $y=C_1e^x+C_2e^{2x}+\frac1{12}e^{5x}$ 是 $y''-3y'+2y=e^{5x}$ 的通解。除代入验证外，还须证明齐次部分线性无关；再求满足 $y(0)=y'(0)=0$ 的特解。
\end{exercise}
\begin{exercise}{798}{三个非齐次解构造通解}
设 $p,q,f$ 在区间 $I$ 上连续，$y_1,y_2,y_3$ 是 $y''+py'+qy=f$ 的三个线性无关的解。证明
\[
y=2C_1y_1+(C_2-C_1)y_2+(1-C_1-C_2)y_3
\]
为通解，并说明两个齐次方向的独立性。再证明此假设蕴含 $f\not\equiv0$。
\end{exercise}
\begin{exercise}{799}{解的叠加与初始条件}
设 $y_1,y_2$ 是方程 $y''+p(x)y'+q(x)y=f(x)$ 的两个解。证明 $\alpha y_1+(1-\alpha)y_2$ 仍为解，并说明一般线性组合成为解所需的条件。取 $f\not\equiv0$ 时，判断 $y_1+y_2$ 与 $y_1-y_2$ 分别满足什么方程。
\end{exercise}
\begin{exercise}{800}{已知解降阶与基本解组}
已知 $y_1=x$ 是 $x^2y''-2xy'+2y=0$ 在 $x>0$ 上的解。令 $y=xv(x)$ 求通解，证明所得两解线性无关，并求满足 $y(1)=0$、$y'(1)=1$ 的解。
\end{exercise}
\subsection{常系数齐次线性微分方程}
\begin{exercise}{801}{候选特解验证与有界解}
分别判断 $e^x$、$e^{2x}$、$1+e^{-2x}$、$xe^{-2x}$ 是否为 $y''+2y'=0$ 的解，给出代入依据。求此方程通解，并确定全实轴上的全部有界解。
\end{exercise}
\begin{exercise}{802}{由解族恢复方程}
求以 $y=C_1e^{6x}+C_2e^{-8x}$ 为通解的首项系数为 $1$ 的二阶常系数齐次方程，证明系数唯一。确定通解中在 $x\to+\infty$ 时趋于零的全部解。
\end{exercise}
\begin{exercise}{803}{共轭特征根与方程恢复}
求以 $e^{2x}\sin3x,e^{2x}\cos3x$ 为基本解组的首项系数为 $1$ 的二阶常系数齐次方程。求满足 $y(0)=0,y'(0)=1$ 的特解，并判断其在 $[0,+\infty)$ 上是否有界。
\end{exercise}
\begin{exercise}{804}{已知指数解与系数信息充分性}
已知 $e^{2x}$ 是 $y''+py'+qy=0$ 的解，求所有可能的 $(p,q)$，说明能否唯一确定系数。若再知 $xe^{2x}$ 也是解，求 $p,q$ 及通解，并解释新增信息与特征根重数的关系。
\end{exercise}
\begin{exercise}{805}{不同实根与初值问题}
求 $y''+5y'+4y=0$ 的通解。再求满足 $y(0)=1,y'(0)=0$ 的解，证明该解在 $x\ge0$ 上恒正且不增加，并求其极限。
\end{exercise}
\begin{exercise}{806}{初值解与极值}
求 $y''+4y'+3y=0$、$y(0)=2,y'(0)=6$ 的特解。确定它在 $[0,+\infty)$ 上的零点及极值，并求最大值。
\end{exercise}
\begin{exercise}{807}{重复特征根与独立解}
求 $4y''-12y'+9y=0$ 的通解，证明所取两个解线性无关。再确定全部满足 $\lim_{x\to+\infty}y(x)=0$ 的解，并说明只消去指数项中的一个系数为何不够。
\end{exercise}
\begin{exercise}{808}{临界阻尼型初值问题}
求 $s''+2s'+s=0$（自变量为 $t$），$s(0)=2,s'(0)=-1$ 的解。证明它在 $t\ge0$ 上恒正且严格减少，求其极限，并说明极限值是否在有限时刻取到。
\end{exercise}
\begin{exercise}{809}{复特征根与零点间隔}
求 $y''+8y'+17y=0$ 的通解。证明任意非零解在 $x\ge0$ 上有无穷多个零点，求相邻零点间距，并判断解在 $x\to+\infty$ 时是否趋于零。
\end{exercise}
\begin{exercise}{810}{衰减振荡与初值确定}
求 $y''+10y'+29y=0$、$y(0)=1,y'(0)=-1$ 的特解。把三角部分化为单一余弦形式，确定其振幅包络，并求第一个正零点。
\end{exercise}
\begin{exercise}{811}{三阶常系数方程与重根}
求 $y'''+y''-y'-y=0$ 的通解，先将特征多项式分组因式分解。证明三个基本解线性无关，并确定全部在 $x\to+\infty$ 时趋于零的解。
\end{exercise}
\begin{exercise}{812}{四阶方程与实复根组合}
求 $y^{(4)}-y'''-y''-y'-2y=0$ 的通解，证明特征多项式可分解为 $(r^2+1)(r-2)(r+1)$。确定通解中在 $[0,+\infty)$ 上有界的解，以及其中趋于零的解。
\end{exercise}
\begin{exercise}{813}{反射导数关系与降阶恢复}
设 $f$ 在实轴上可导，满足 $f'(x)=f(-x)$、$f(0)=1$。证明 $f$ 二阶可导并求 $f$，再代回反射关系核验，说明只满足导出的二阶方程为何还不足以完成验证。
\end{exercise}
\begin{exercise}{814}{两个微分关系的相容性}
设 $f$ 同时满足 $f''+f'-2f=0$ 与 $f''+f=2e^x$。求 $f$，先分别给出两个方程解的形式，再确定公共解，并证明其唯一性。
\end{exercise}
\begin{exercise}{815}{重根、参数与长期行为}
求初值问题 $y''+2ay'+a^2y=0$、$y(0)=0$、$y'(0)=1$ 的解。讨论实参数 $a$ 使解在 $[0,+\infty)$ 上有界或趋于零的条件，并在 $a>0$ 时求其最大值。
\end{exercise}
\begin{exercise}{816}{阻尼分类与无零点条件}
考虑 $y''+2\zeta y'+y=0$，$y(0)=1$、$y'(0)=0$，其中 $\zeta\ge0$。分别求 $0\le\zeta<1$、$\zeta=1$、$\zeta>1$ 时的解，并确定解在 $x>0$ 上始终为正的参数范围。
\end{exercise}
\subsection{常系数非齐次线性微分方程}
\begin{exercise}{817}{多项式特解与通解}
求 $y''-4y'+3y=x$ 的一个一次多项式特解及通解。证明一次多项式特解唯一，并求满足 $y(0)=y'(0)=0$ 的特解。
\end{exercise}
\begin{exercise}{818}{多项式外力与长期行为}
求 $y''+5y'+4y=3-2x$ 的通解。证明所有解与同一个一次多项式之差在 $x\to+\infty$ 时趋于零，并写出该多项式。
\end{exercise}
\begin{exercise}{819}{零特征根与待定系数形式}
求 $2y''+5y'=5x^2-2x-1$ 的通解。先说明特解为何需要三次多项式；再求满足 $y(0)=y'(0)=0$ 的特解，并核验常数项可并入齐次解。
\end{exercise}
\begin{exercise}{820}{由齐次通解恢复系数与求特解}
已知 $y''+ay'+by=0$ 的通解为 $(C_1+C_2x)e^x$。求 $a,b$，再求 $y''+ay'+by=x$ 满足 $y(0)=2,y'(0)=0$ 的解，分别核验方程及初值。
\end{exercise}
\begin{exercise}{821}{混合外力与共振特解形式}
对 $y''-5y'+6y=3x+2e^{2x}$，说明为何可设特解为 $Ax+B+Cxe^{2x}$，求出系数并写出通解。解释指数项前乘 $x$ 的原因，以及单设 $Ce^{2x}$ 为什么无效。
\end{exercise}
\begin{exercise}{822}{指数多项式与单根共振}
求 $y''-3y'+2y=xe^{2x}$ 的通解。先写出最小的待定系数特解形式，再求满足 $y(0)=y'(0)=0$ 的特解，并用代换 $y=e^{2x}v$ 核验。
\end{exercise}
\begin{exercise}{823}{已知非齐次解反求常系数}
已知 $y=e^{-2x}+(x^2+2)e^x$ 是 $y''+ay'+by=(cx+d)e^x$ 的解，其中 $a,b,c,d$ 为常数。求全部系数并证明其唯一性，再写出该方程通解。
\end{exercise}
\begin{exercise}{824}{三个解恢复方程与外力}
某首项系数为 $1$ 的二阶常系数非齐次方程有三个解 $y_1=e^{2x}+xe^x$、$y_2=e^{-x}+xe^x$、$y_3=e^{2x}-e^{-x}+xe^x$。求方程及通解，并说明如何先从解的差恢复齐次部分，再恢复外力项。
\end{exercise}
\begin{exercise}{825}{复根共振与候选解验证}
对 $y''-2y'+5y=e^x\cos2x$，分别检验 $xe^x(\cos2x+\sin2x)$、$e^x(\cos2x+\sin2x)$、$\frac14xe^x\cos2x$、$\frac14xe^x\sin2x$ 是否为特解。选出成立者，写出通解，并说明共振所需的 $x$ 因子。
\end{exercise}
\begin{exercise}{826}{倍角分解与三角共振}
求 $y''+4y=\cos^2x$ 的一个特解及通解。先把外力分解为常数项与三角项，分别说明其待定系数形式；再求满足 $y(0)=y'(0)=0$ 的解。
\end{exercise}
\begin{exercise}{827}{多项式三角外力}
求 $y''+4y=x\cos x$ 的通解。说明仅设特解为 $Ax\cos x$ 是否足够，并写出包含必要正弦项的特解形式；再求满足 $y(0)=y'(0)=0$ 的特解。
\end{exercise}
\begin{exercise}{828}{外力叠加与指数共振}
求 $y''-y=e^x+4\cos x$ 的通解。分别处理两种外力并说明指数项的共振，再求满足 $y(0)=y'(0)=0$ 的特解，判断其在 $[0,+\infty)$ 上是否有界。
\end{exercise}
\begin{exercise}{829}{共振与参数分类}
对 $\omega>0$，求初值问题 $y''+y=\cos\omega x$、$y(0)=y'(0)=0$ 的解，分别讨论 $\omega\ne1$ 与 $\omega=1$。判断解有界的参数条件，并说明共振时无界项的来源。
\end{exercise}
\begin{exercise}{830}{指数多项式外力与特解形式}
求 $y''-2y'+y=e^x(x+1)$ 的通解及满足 $y(0)=y'(0)=0$ 的特解。先写出待定系数法所需的特解形式，并用 $y=e^xv$ 的代换核验计算结果。
\end{exercise}
